
\section*{Part III: A Starting Strategy for Bayes Institute}
\addcontentsline{toc}{section}{Part III: A Starting Strategy for Bayes Institute}

\subsection{14. The actual initial position}

Bayes does not begin with 0 \textit{personal} capability. Its founder has engineering, AI/\allowbreak ML, statistics, teaching, curriculum design, and industry experience. But organizationally Bayes begins close to 0: no validated enduring product, limited cash available to risk, no alumni body, no mature faculty system, no degree authority, and no guaranteed distribution channel.

The founder also intends to retain a full-time engineering job while running a lean operation. That constraint is not embarrassing. It is the economic basis for avoiding desperate sales or outside capital on unfavorable terms. It does, however, make founder time a scarce and valuable resource. Early proposals must fit a sustainable operating envelope of roughly 15-25 hours per week, not an imaginary 70-hour 2nd job.

The founding mission - \textit{to create the conditions in which human possibility becomes human capability} - should be treated as stable. The 1st course topic, pricing, audience, cadence, and sales channel should not.

\subsection{15. 3 initial engines: why they are not interchangeable}

\textbf{Engine A: D2C technical cohorts.} Potentially strong control over learner experience, brand, assessment, and alumni relationships. Cash may come before delivery, but the founder must build trust and acquire learners 1 by 1. Failure modes: high CAC, weak completion, dependence on a star teacher.

\textbf{Engine B: College partnerships.} Potential for larger contracts and wide reach; credible access to learners, facilities, and existing academic pathways. Failure modes: procurement delays, low payments, slow collections, institutional politics, faculty resistance, and being reduced to a commodity placement vendor.

\textbf{Engine C: Enterprise training and applied capability building.} Can deliver significant revenue per customer, access valuable real problems, and create research opportunities. Failure modes: custom consulting traps, dependence on a handful of clients, long sales cycles, restrictive IP, confidentiality obligations, and service-delivery intensity.

These are \textbf{separate market hypotheses}. They share expertise but not necessarily buyers, budgets, or operational rhythms. Avoid announcing all 3 as fully functioning divisions. Test them as options, then choose the 1 with the clearest near-term proof and strongest accumulation of future assets.

\subsection{16. Why a 1st education product must be narrower than the mission}

``Learn AI'' is a topic. ``Become more'' is an institutional philosophy. Neither is a sufficiently narrow 1st purchasing proposition.

A useful 1st product identifies:

\begin{enumerate}[leftmargin=*,itemsep=1.5pt]
  \item A buyer with a pressing consequence of inaction.
  \item A baseline capability that can be measured.
  \item A destination capability whose evidence can be inspected.
  \item A realistic duration and weekly effort.
  \item A delivery method that works for the founder's time budget.
  \item A price the buyer will actually commit to.
\end{enumerate}

For example, \textbf{\textbf{[Hypothesis]}} a selective program for working software engineers who need to design, evaluate, secure, and deploy reliable AI applications. A learner should leave able to justify architecture, write meaningful tests, measure failure modes, and defend design tradeoffs on a new problem. This is a possible wedge, not a validated recommendation.

The early Bayesian identity should be built from high standards, credible teaching, and authentic work - not synthetic exclusivity or claims of guaranteed mobility.

\subsection{17. Design the 1st 90 days as experiments, not a launch campaign}

\textbf{Days 1-15:} Interview 20-30 genuinely relevant prospective learners and 5-10 employers or hiring managers about costly capability failures. Ask for decisions and evidence, not abstract excitement. Document the language of their problem, their alternatives, and whether they currently spend money to solve it.

\textbf{Days 16-30:} Publish a concrete competency rubric and 2 demonstrations of expert teaching. Offer an application and a paid diagnostic/\allowbreak workshop with a narrowly defined deliverable. Observe who commits funds, not only who signs up.

\textbf{Days 31-60:} Deliver a paid, bounded pilot to roughly 8-15 learners if demand supports it. Use blind or independently graded baseline and post-program tasks. Collect the full cost of delivery, including founder time imputed at a reasonable replacement rate.

\textbf{Days 61-90:} Review demand, learning, price, time, referrals, and delivery problems. Choose 1 of 3 actions: repeat, revise, or stop. Do not scale weak learning evidence with stronger advertising.

If a college or enterprise signs a serious paid pilot sooner than learners do, test that path. But require a small defined contract, acceptance criteria, payment terms, and a repeatability hypothesis. Do not build 3 organizations in parallel because 3 types of prospects sounded interested.

\subsection{18. The 1st role is not necessarily an instructor}

With a working founder, the bottleneck may initially be scheduling, learner support, lead follow-up, feedback logistics, payment collection, video production, or grading. Hire the task that returns the most \textit{high-value founder hours} while preserving quality. A part-time learning operations specialist may be worth more than a 2nd expensive star instructor at 12 learners.

A candidate role should be justified by this equation:

\textbf{Net value of hire = contribution from additional capacity + founder hours recovered $\times$ economic value of those hours + quality/\allowbreak risk improvements - fully loaded cost - management burden.}

The recovered hours are real only if the founder actually uses them for sales, teaching design, hiring, or rest - not if administrative work simply re-expands.

\subsection{19. Early service revenue: fuel, teacher, or prison?}

Services can be financially attractive, especially if a founder already knows how to sell and deliver them. But they come in 3 forms:

\begin{itemize}[leftmargin=*,itemsep=1.5pt]
  \item \textbf{Mission-aligned service:} an enterprise AI evaluation lab that creates reusable case studies and relationships.
  \item \textbf{Cash-funding service:} a well-paid short contract that does little for the long-term educational vision but funds reserves.
  \item \textbf{Strategic trap:} repeated custom projects that require the founder continuously and crowd out the institutional product.
\end{itemize}

All 3 may generate profit. The 1st has compounding institutional value; the 2nd can be acceptable within explicit limits; the 3rd is a threat unless it can be delegated or priced very highly.

\textbf{\textbf{[Decision rule]}} For a side-venture founder, put a hard cap on custom projects: time, revenue concentration, delivery commitments, and the maximum weeks they can delay the core learning product. Never build a brand that promises a research university while functionally operating as an outsourced engineering agency without saying so.

-


